\documentclass[11pt,a4paper]{article}
\usepackage[margin=1in]{geometry}
\usepackage{amsmath,amssymb,amsthm}
\usepackage{algorithm}
\usepackage{algpseudocode}
\usepackage{booktabs}
\usepackage{hyperref}

\theoremstyle{definition}
\newtheorem{definition}{Definition}[section]
\newtheorem{theorem}{Theorem}[section]
\newtheorem{lemma}[theorem]{Lemma}
\newtheorem{remark}{Remark}[section]

\title{\textbf{Rigorous Mathematical Specification, Decoupled Weight Decay Mechanics, and Convergence Analysis of AdamW}}
\author{Team Scaibu \\ \small Email: \texttt{jaydeep.9664920749@gmail.com} \\ \small Architecture and Core Engine Team}
\date{October 2026}

\begin{document}
\maketitle

\begin{abstract}
This document presents the complete theoretical and algorithmic specification of Decoupled Weight Decay Adam (AdamW). In standard adaptive gradient methods, $L_2$ regularization is incorrectly coupled into the moving average moment buffers, causing weights with large historical gradients to be regularized less than weights with small gradients. AdamW restores true $L_2$ weight decay by applying the shrinkage operation directly to the parameter vector independently of moment accumulation. We formalize this distinction, provide the exact algorithm pseudo-code matching the reference implementation, prove regularization equivalence, analyze algorithmic complexity, compute a step-by-step numerical example, and discuss floating-point numerical considerations.
\end{abstract}

\section{Notation and Mathematical Preliminaries}

Let $\theta \in \mathbb{R}^P$ denote the network parameter vector optimized across discrete training steps $t \in \{1, 2, \dots, T\}$ over loss function $\mathcal{L}(\theta)$.

\begin{itemize}
    \item $g_t \triangleq \nabla_\theta \mathcal{L}(\theta_t) \in \mathbb{R}^P$: Pure loss objective gradient (excluding weight decay).
    \item $m_t \in \mathbb{R}^P$: Exponential moving average (first moment) of pure loss gradients.
    \item $v_t \in \mathbb{R}^P$: Exponential moving average (second moment) of squared pure loss gradients.
    \item $\hat{m}_t, \hat{v}_t \in \mathbb{R}^P$: Bias-corrected moment estimates.
    \item $\lambda \ge 0$: Decoupled weight decay coefficient.
    \item $\eta > 0$: Base learning rate (step size).
    \item $\beta_1, \beta_2 \in [0, 1)$: Exponential decay hyperparameters for moment buffers.
    \item $\epsilon > 0$: Regularization constant ensuring division safety.
\end{itemize}

\begin{table}[h]
\centering
\caption{Summary of Mathematical Symbols for AdamW}
\begin{tabular}{lll}
\toprule
\textbf{Symbol} & \textbf{Type / Domain} & \textbf{Description} \\
\midrule
$\theta_t$ & $\mathbb{R}^P$ & Parameter iterate at step $t$ \\
$g_t$ & $\mathbb{R}^P$ & Gradient of pure unregularized loss $\nabla \mathcal{L}(\theta_t)$ \\
$m_t, v_t$ & $\mathbb{R}^P$ & First and second raw moment buffers \\
$\hat{m}_t, \hat{v}_t$ & $\mathbb{R}^P$ & Bias-corrected moment vectors \\
$\lambda$ & $\mathbb{R}_{\ge 0}$ & Decoupled weight decay rate ($0.01$) \\
$\beta_1, \beta_2$ & $[0, 1)$ & Moment decay parameters ($0.9, 0.999$) \\
$\eta$ & $\mathbb{R}_{> 0}$ & Base learning rate ($10^{-3}$) \\
$\epsilon$ & $\mathbb{R}_{> 0}$ & Stability constant ($10^{-8}$) \\
\bottomrule
\end{tabular}
\end{table}

\section{Problem Definition and Core Concepts}

\subsection{Intuition and Motivation: $L_2$ Regularization vs. True Weight Decay}
In standard SGD, adding an $L_2$ penalty $\frac{\lambda}{2} \|\theta\|_2^2$ to the loss yields gradient $\tilde{g}_t = g_t + \lambda \theta_t$, producing update:
\begin{equation}
\theta_{t+1} = \theta_t - \eta (g_t + \lambda \theta_t) = (1 - \eta \lambda) \theta_t - \eta g_t
\end{equation}
which is mathematically identical to decaying weights by $(1 - \eta \lambda)$.
However, in standard Adam with $L_2$ regularization, $\tilde{g}_t = g_t + \lambda \theta_t$ is passed directly into moment updates $m_t$ and $v_t$. Consequently, the effective shrinkage applied to parameter $\theta_i$ becomes:
\begin{equation}
\Delta \theta_i \approx -\frac{\eta \lambda}{\sqrt{\hat{v}_i} + \epsilon} \theta_i
\end{equation}
Thus, parameters with large historical gradients (large $\hat{v}_i$) experience \textit{smaller} weight decay, while parameters with infrequent or small gradients experience \textit{disproportionately larger} weight decay. AdamW resolves this fundamental pathology by decoupling the weight decay step from gradient moment accumulation.

\subsection{Formal Mathematical Definition}
\begin{definition}[AdamW Update Rule]
For step $t \ge 1$, given parameters $\theta_{t-1}$, pure gradient $g_t$, moments $m_{t-1}, v_{t-1}$, and hyperparameters $(\eta, \lambda, \beta_1, \beta_2, \epsilon)$:
\begin{align}
\theta_t^{\text{decay}} &= \theta_{t-1} \odot (1 - \eta \lambda) \label{eq:adamw_decay} \\
m_t &= \beta_1 m_{t-1} + (1 - \beta_1) g_t \label{eq:adamw_m} \\
v_t &= \beta_2 v_{t-1} + (1 - \beta_2) (g_t \odot g_t) \label{eq:adamw_v} \\
\hat{m}_t &= \frac{m_t}{1 - \beta_1^t}, \quad \hat{v}_t = \frac{v_t}{1 - \beta_2^t} \label{eq:adamw_bc} \\
\theta_t &= \theta_t^{\text{decay}} - \eta \frac{\hat{m}_t}{\sqrt{\hat{v}_t} + \epsilon} \label{eq:adamw_update}
\end{align}
\end{definition}

\section{Algorithmic Specification}

The exact computational execution implemented in \texttt{NnAlgoAdamwOptimizer} is detailed below.

\begin{algorithm}[H]
\caption{Decoupled Weight Decay Adam (AdamW) Step Update}
\begin{algorithmic}[1]
\Require Parameters $\theta \in \mathbb{R}^P$, gradients $g \in \mathbb{R}^P$, moment buffers $m, v \in \mathbb{R}^P$
\Require Step index $t \ge 1$, learning rate $\eta > 0$, weight decay $\lambda \ge 0$, decays $\beta_1, \beta_2 \in [0, 1)$, epsilon $\epsilon > 0$
\Ensure Updated parameters $\theta_{\text{next}}$, updated moments $m_{\text{next}}, v_{\text{next}} \in \mathbb{R}^P$
\State $P \gets \text{length}(\theta)$
\If{$P = 0$ \textbf{or} $\text{length}(g) \ne P$ \textbf{or} $\text{length}(m) \ne P$ \textbf{or} $\text{length}(v) \ne P$}
    \State \textbf{throw} PreconditionFailureException(``Buffer dimensions mismatch or empty'')
\EndIf
\If{$t < 1$ \textbf{or} $\eta \le 0$ \textbf{or} $\epsilon \le 0$ \textbf{or} $\lambda < 0$ \textbf{or} $\beta_1 < 0$ \textbf{or} $\beta_1 \ge 1$ \textbf{or} $\beta_2 < 0$ \textbf{or} $\beta_2 \ge 1$}
    \State \textbf{throw} PreconditionFailureException(``Hyperparameter out of valid range'')
\EndIf
\State $b_1 \gets 1.0 - \beta_1^t$
\State $b_2 \gets 1.0 - \beta_2^t$
\State $\theta_{\text{next}} \gets \text{new Array of size } P$
\State $m_{\text{next}} \gets \text{new Array of size } P$
\State $v_{\text{next}} \gets \text{new Array of size } P$
\For{$i \gets 1 \textbf{ to } P$}
    \State $\theta_{\text{decayed}} \gets \theta[i] \cdot (1.0 - \eta \cdot \lambda)$
    \State $m_{\text{next}}[i] \gets \beta_1 \cdot m[i] + (1.0 - \beta_1) \cdot g[i]$
    \State $v_{\text{next}}[i] \gets \beta_2 \cdot v[i] + (1.0 - \beta_2) \cdot (g[i])^2$
    \State $\hat{m} \gets m_{\text{next}}[i] / b_1$
    \State $\hat{v} \gets v_{\text{next}}[i] / b_2$
    \State $\theta_{\text{next}}[i] \gets \theta_{\text{decayed}} - \left(\frac{\eta \cdot \hat{m}}{\sqrt{\hat{v}} + \epsilon}\right)$
\EndFor
\State \Return $\{\text{updated\_parameters}: \theta_{\text{next}}, \text{updated\_exp\_avg}: m_{\text{next}}, \text{updated\_exp\_avg\_sq}: v_{\text{next}}\}$
\end{algorithmic}
\end{algorithm}

\section{Theoretical Guarantees and Invariance Analysis}

\begin{theorem}[Decoupled Regularization Invariance]
Under AdamW, the rate of relative weight decay per step $\frac{\|\theta_t^{\text{decay}} - \theta_{t-1}\|_2}{\|\theta_{t-1}\|_2} = \eta \lambda$ is strictly uniform across all coordinates and invariant to gradient scale and moment variance.
\end{theorem}

\begin{proof}
By Equation~\eqref{eq:adamw_decay}, for every coordinate $i \in \{1, \dots, P\}$:
\begin{equation}
\theta_{t, i}^{\text{decay}} - \theta_{t-1, i} = \theta_{t-1, i}(1 - \eta \lambda) - \theta_{t-1, i} = -\eta \lambda \theta_{t-1, i}
\end{equation}
Computing the Euclidean norm across all $P$ coordinates:
\begin{equation}
\|\theta_t^{\text{decay}} - \theta_{t-1}\|_2 = \sqrt{\sum_{i=1}^P (-\eta \lambda \theta_{t-1, i})^2} = \eta \lambda \sqrt{\sum_{i=1}^P \theta_{t-1, i}^2} = \eta \lambda \|\theta_{t-1}\|_2
\end{equation}
Dividing by $\|\theta_{t-1}\|_2$ yields exactly $\eta \lambda$, independent of $g_{t, i}$, $m_{t, i}$, or $v_{t, i}$.
\end{proof}

\subsection{Limitations and Best Practices}
\begin{itemize}
    \item \textbf{LayerNorm / Bias Exclusions}: In Transformers and residual networks, decoupled weight decay should typically be disabled ($\lambda = 0$) for 1D parameter vectors (biases, normalization scale/shift gains).
    \item \textbf{Learning Rate Schedule Coupling}: Because shrinkage is $\eta_t \lambda$, decaying the learning rate $\eta_t$ automatically decreases the weight decay rate proportionally.
\end{itemize}

\section{Complexity Analysis}

\subsection{Time Complexity}
\begin{itemize}
    \item \textbf{Decay Step}: Exactly 1 subtraction and 1 multiplication per parameter ($\mathcal{O}(P)$).
    \item \textbf{Moment and Update Steps}: Standard Adam FLOPs ($\mathcal{O}(P)$).
    \item \textbf{Total Time}: $\Theta(P)$ FLOPs per step.
\end{itemize}

\subsection{Space Complexity}
\begin{itemize}
    \item \textbf{Persistent State}: Stores $m \in \mathbb{R}^P$ and $v \in \mathbb{R}^P$.
    \item \textbf{Auxiliary Memory}: Exactly $2P$ floating point elements ($\Theta(P)$ auxiliary space).
\end{itemize}

\section{Worked Numerical Example}

Let $P=1$, $\theta = [2.0]$, $g = [0.1]$, $m = [0.0]$, $v = [0.0]$, step $t=1$, $\eta = 0.01$, $\lambda = 0.05$, $\beta_1 = 0.9$, $\beta_2 = 0.999$, $\epsilon = 10^{-8}$.

\begin{enumerate}
    \item \textbf{Decay computation}:
    \begin{equation}
    \theta_{\text{decayed}} = 2.0 \times (1.0 - 0.01 \times 0.05) = 2.0 \times (1.0 - 0.0005) = 2.0 \times 0.9995 = 1.9990
    \end{equation}
    \item \textbf{Moment updates}:
    \begin{align}
    m_{\text{next}} &= 0.9(0) + 0.1(0.1) = 0.0100 \\
    v_{\text{next}} &= 0.999(0) + 0.001(0.1^2) = 0.0000100
    \end{align}
    \item \textbf{Bias corrections}:
    \begin{align}
    \hat{m} &= \frac{0.01}{1 - 0.9^1} = \frac{0.01}{0.1} = 0.1000 \\
    \hat{v} &= \frac{0.00001}{1 - 0.999^1} = \frac{0.00001}{0.001} = 0.0100
    \end{align}
    \item \textbf{Adaptive Step}:
    \begin{equation}
    \Delta \theta = \frac{0.01 \times 0.1000}{\sqrt{0.01} + 10^{-8}} = \frac{0.0010}{0.10000001} \approx 0.009999999
    \end{equation}
    \item \textbf{Final parameter}:
    \begin{equation}
    \theta_{\text{next}} = 1.9990 - 0.009999999 = 1.989000001
    \end{equation}
\end{enumerate}

\section{Numerical Considerations and Edge Cases}

\begin{itemize}
    \item \textbf{Order of Operations}: Performing $\theta \leftarrow \theta(1 - \eta \lambda)$ before the adaptive gradient step maintains numerical precision and avoids catastrophic cancellation.
    \item \textbf{Multi-Precision Mixed FP16/BF16 Training}: Master weights must be stored and updated in FP32, as the decay factor $\eta \lambda \approx 10^{-5}$ is below the minimum representable precision delta in FP16.
\end{itemize}

\begin{thebibliography}{9}
\bibitem{loshchilov2019} I.~Loshchilov and F.~Hutter, ``Decoupled weight decay regularization,'' in \emph{International Conference on Learning Representations (ICLR)}, 2019.
\bibitem{kingma2015} D.~P.~Kingma and J.~Ba, ``Adam: A method for stochastic optimization,'' in \emph{International Conference on Learning Representations (ICLR)}, 2015.
\bibitem{touvron2023} H.~Touvron et al., ``Llama: Open and efficient foundation language models,'' \emph{arXiv preprint arXiv:2302.13971}, 2023.
\end{thebibliography}

\end{document}
